% GENERATED FILE. Edit canonical lessons, then run npm run generate:books.
% canonical-source-hash: fnv1a64:4f05f55af14ccbb1
% canonical-lessons: AR-C402-hatman, AR-C402-bi-d-darura, AR-C402-nisbiyyan, AR-C402-itlaqan, AR-C402-yaqin

\chapter{How sure}
\label{ch:ar-how-sure}
\hlchaptermodality{402}

\begin{chapteropening}
\textbf{By the end of this chapter:} \emph{I can say inevitably, necessarily, relatively and not at all, and talk about certainty.}

Complete the last lesson of chapter 402: I can say inevitably, necessarily, relatively and not at all, and talk about certainty.
\end{chapteropening}

\section[\texorpdfstring{\ar{حتما} (ḥatman)}{حتما (ḥatman)}]{\ar{حتما} (ḥatman) — inevitably}

\label{lesson:AR-C402-hatman}



\begin{quote}
Before the new one: say the Arabic for to consider likely, then the Arabic for to rule out.
\end{quote}

\subsection*{You'll want to know: \ar{حتما}}
\textbf{\ar{حتما}} — \emph{ḥatman} — \textquotedblleft{}inevitably\textquotedblright{}.

\begin{quote}
\ar{سينجح} \ar{حتما}.
\end{quote}

\begin{quote}
\emph{sa-yanjaḥu ḥatman.}
\end{quote}

\begin{quote}
He is bound to succeed.
\end{quote}

\begin{cousinweb}
\textbf{\ar{حتما}} (\emph{ḥatman}) is built on the native Arabic root \emph{ḥ-t-m}, \textquotedblleft{}to make firm, to decree\textquotedblright{} (Lane). \emph{Ḥatman} is the verbal noun \emph{ḥatm}, decreeing, with the adverb ending \emph{-an}: as decreed, inevitably.
\end{cousinweb}

\begin{grammarlens}[title={What you've built}]
One more word for saying what happens if.
\end{grammarlens}

\subsection*{Your turn}
\begin{itemize}
\raggedright
  \item \emph{Say it:} \emph{ḥatman}
  \item \emph{Say it:} \emph{ḥatman}, once more
  \item \emph{Say it:} \emph{sa-yanjaḥu ḥatman.} — He is bound to succeed.
  \item \emph{Recall:} say the Arabic for to consider likely, then the Arabic for to rule out, then say \emph{ḥatman} again
\end{itemize}

\begin{tcolorbox}[breakable,colback=teal!4,colframe=teal!35!black,title={Before you move on}]
What does \ar{حتما} mean? (\textquotedblleft{}Inevitably\textquotedblright{}.) Say it once more.
\end{tcolorbox}
\section[\texorpdfstring{\ar{بالضرورة} (bi-ḍ-ḍarūra)}{بالضرورة (bi-ḍ-ḍarūra)}]{\ar{بالضرورة} (bi-ḍ-ḍarūra) — necessarily}

\label{lesson:AR-C402-bi-d-darura}



\begin{quote}
Before the new one: say the Arabic for to rule out, then the Arabic for inevitably.
\end{quote}

\subsection*{You'll want to know: \ar{بالضرورة}}
\textbf{\ar{بالضرورة}} — \emph{bi-ḍ-ḍarūra} — \textquotedblleft{}necessarily\textquotedblright{}.

\begin{quote}
\ar{الغالي} \ar{ليس} \ar{بالضرورة} \ar{أفضل}.
\end{quote}

\begin{quote}
\emph{al-ghālī laysa bi-ḍ-ḍarūra afḍal.}
\end{quote}

\begin{quote}
Expensive is not necessarily better.
\end{quote}

\begin{cousinweb}
\textbf{\ar{بالضرورة}} (\emph{bi-ḍ-ḍarūra}) is built on the native Arabic root \emph{ḍ-r-r}, \textquotedblleft{}to harm; to compel\textquotedblright{} (Lane). \emph{Ḍarūra} is necessity, what one is compelled to; \emph{bi-ḍ-ḍarūra} is by necessity. \ar{ضروري}, necessary, is from the same root.
\end{cousinweb}

\begin{grammarlens}[title={What you've built}]
One more word for saying what happens if.
\end{grammarlens}

\subsection*{Your turn}
\begin{itemize}
\raggedright
  \item \emph{Say it:} \emph{bi-ḍ-ḍarūra}
  \item \emph{Say it:} \emph{bi-ḍ-ḍarūra}, once more
  \item \emph{Say it:} \emph{al-ghālī laysa bi-ḍ-ḍarūra afḍal.} — Expensive is not necessarily better.
  \item \emph{Recall:} say the Arabic for to rule out, then the Arabic for inevitably, then say \emph{bi-ḍ-ḍarūra} again
\end{itemize}

\begin{tcolorbox}[breakable,colback=teal!4,colframe=teal!35!black,title={Before you move on}]
What does \ar{بالضرورة} mean? (\textquotedblleft{}Necessarily\textquotedblright{}.) Say it once more.
\end{tcolorbox}
\section[\texorpdfstring{\ar{نسبيا} (nisbiyyan)}{نسبيا (nisbiyyan)}]{\ar{نسبيا} (nisbiyyan) — relatively}

\label{lesson:AR-C402-nisbiyyan}



\begin{quote}
Before the new one: say the Arabic for inevitably, then the Arabic for necessarily.
\end{quote}

\subsection*{You'll want to know: \ar{نسبيا}}
\textbf{\ar{نسبيا}} — \emph{nisbiyyan} — \textquotedblleft{}relatively\textquotedblright{}.

\begin{quote}
\ar{الشقة} \ar{رخيصة} \ar{نسبيا}.
\end{quote}

\begin{quote}
\emph{ash-shaqqa rakhīṣa nisbiyyan.}
\end{quote}

\begin{quote}
The flat is relatively cheap.
\end{quote}

\begin{cousinweb}
\textbf{\ar{نسبيا}} (\emph{nisbiyyan}) is built on the native Arabic root \emph{n-s-b}, \textquotedblleft{}to trace lineage, to relate\textquotedblright{} (Lane). \emph{Nisbiyyan} is the adverb of \emph{nisbī}, relative, from a noun meaning a relation.
\end{cousinweb}

\begin{grammarlens}[title={What you've built}]
One more word for saying what happens if.
\end{grammarlens}

\subsection*{Your turn}
\begin{itemize}
\raggedright
  \item \emph{Say it:} \emph{nisbiyyan}
  \item \emph{Say it:} \emph{nisbiyyan}, once more
  \item \emph{Say it:} \emph{ash-shaqqa rakhīṣa nisbiyyan.} — The flat is relatively cheap.
  \item \emph{Recall:} say the Arabic for inevitably, then the Arabic for necessarily, then say \emph{nisbiyyan} again
\end{itemize}

\begin{tcolorbox}[breakable,colback=teal!4,colframe=teal!35!black,title={Before you move on}]
What does \ar{نسبيا} mean? (\textquotedblleft{}Relatively\textquotedblright{}.) Say it once more.
\end{tcolorbox}
\section[\texorpdfstring{\ar{إطلاقا} (iṭlāqan)}{إطلاقا (iṭlāqan)}]{\ar{إطلاقا} (iṭlāqan) — not at all}

\label{lesson:AR-C402-itlaqan}



\begin{quote}
Before the new one: say the Arabic for necessarily, then the Arabic for relatively.
\end{quote}

\subsection*{You'll want to know: \ar{إطلاقا}}
\textbf{\ar{إطلاقا}} — \emph{iṭlāqan} — \textquotedblleft{}not at all\textquotedblright{}.

\begin{quote}
\ar{لا} \ar{أحب} \ar{القهوة} \ar{إطلاقا}.
\end{quote}

\begin{quote}
\emph{lā uḥibbu al-qahwa iṭlāqan.}
\end{quote}

\begin{quote}
I do not like coffee at all.
\end{quote}

With a negative: \ar{لا} … \ar{إطلاقا}, not … at all.

\begin{cousinweb}
\textbf{\ar{إطلاقا}} (\emph{iṭlāqan}) is built on the native Arabic root \emph{ṭ-l-q}, \textquotedblleft{}to be free, let loose\textquotedblright{} (Lane). \emph{Iṭlāqan} is the Form IV verbal noun, setting loose, with the adverb ending \emph{-an}: without limit, and with a negative, not at all.
\end{cousinweb}

\begin{grammarlens}[title={What you've built}]
One more word for saying what happens if.
\end{grammarlens}

\subsection*{Your turn}
\begin{itemize}
\raggedright
  \item \emph{Say it:} \emph{iṭlāqan}
  \item \emph{Say it:} \emph{iṭlāqan}, once more
  \item \emph{Say it:} \emph{lā uḥibbu al-qahwa iṭlāqan.} — I do not like coffee at all.
  \item \emph{Recall:} say the Arabic for necessarily, then the Arabic for relatively, then say \emph{iṭlāqan} again
\end{itemize}

\begin{tcolorbox}[breakable,colback=teal!4,colframe=teal!35!black,title={Before you move on}]
What does \ar{إطلاقا} mean? (\textquotedblleft{}Not at all\textquotedblright{}.) Say it once more.
\end{tcolorbox}
\section[\texorpdfstring{\ar{يقين} (yaqīn)}{يقين (yaqīn)}]{\ar{يقين} (yaqīn) — certainty}

\label{lesson:AR-C402-yaqin}



\begin{quote}
Before the new one: say the Arabic for relatively, then the Arabic for not at all.
\end{quote}

\subsection*{You'll want to know: \ar{يقين}}
\textbf{\ar{يقين}} — \emph{yaqīn} — \textquotedblleft{}certainty\textquotedblright{}.

\begin{quote}
\ar{عندي} \ar{يقين} \ar{بأنه} \ar{صادق}.
\end{quote}

\begin{quote}
\emph{ʿindī yaqīn bi-annahu ṣādiq.}
\end{quote}

\begin{quote}
I am certain that he is honest.
\end{quote}

\begin{cousinweb}
\textbf{\ar{يقين}} (\emph{yaqīn}) is built on the native Arabic root \emph{y-q-n}, usually traced to the sense \textquotedblleft{}to be sure\textquotedblright{}. \emph{Yaqīn} is certainty, knowledge without doubt.
\end{cousinweb}

\begin{grammarlens}[title={What you've built}]
One more word for saying what happens if.
\end{grammarlens}

\subsection*{Your turn}
\begin{itemize}
\raggedright
  \item \emph{Say it:} \emph{yaqīn}
  \item \emph{Say it:} \emph{yaqīn}, once more
  \item \emph{Say it:} \emph{ʿindī yaqīn bi-annahu ṣādiq.} — I am certain that he is honest.
  \item \emph{Recall:} say the Arabic for relatively, then the Arabic for not at all, then say \emph{yaqīn} again
\end{itemize}

\begin{tcolorbox}[breakable,colback=teal!4,colframe=teal!35!black,title={Before you move on}]
What does \ar{يقين} mean? (\textquotedblleft{}Certainty\textquotedblright{}.) Say it once more.
\end{tcolorbox}
